\sisection{S1. Data, estimands and overlap}

\sisubsection{S1.1 Row-level sources and preprocessing}

Two docking score matrices carry every spectral result in the paper. Both are secondary use
of released tables. We ran no docking and made no measurement.

Docking-44 is the frozen export of Nikitin et al.~\cite{nikitin2025}, redistributed as
\path{data/frozen/df_final_v4.csv.gz} under CC BY 4.0: 13{,}283{,}042 bytes, SHA-256
beginning \texttt{25436e49b9fc}. The carrier file has 47 PDB-named score columns. Every analysis
selects the 44 structures mapped to the 44 toxicity-relevant targets in the published source
panel; the appended 8pm6 (BSEP), 6itm (FXR) and 9fzj (PXR) columns are outside that panel and are
excluded before score statistics are computed. The selected matrix therefore holds 12{,}651
ligands by 44 Vina-GPU~2.0 target columns \cite{ding2023}, or 556{,}644 cells. Of these, 8{,}159 cells are missing, which is 1.4657\%
of the matrix, and they sit in 3{,}354 ligand rows. One target, 4mqs, holds 3{,}352 of the
8{,}159 and is missing for 26.50\% of rows; the next largest, 8st0, holds 1{,}007. We
replaced each missing cell by the mean of its target column. No cell in the selected 44-column
matrix is strictly positive, so the clip-at-zero rule removed nothing, and we keep it because it is the
released preprocessing contract.

DOCKSTRING is DOCKSTRING v1 \cite{dockstring2022}, redistributed as
\path{data/frozen/dockstring-dataset.tsv.gz} under Apache-2.0: 26{,}561{,}332 bytes, SHA-256
beginning \texttt{e15a58258dbd}. The source table has 260{,}155 rows by 58 AutoDock Vina
columns \cite{vina2010}. Ninety-five rows carry at least one of the 339 missing cells, and we
dropped those rows rather than impute them, leaving 260{,}060 complete rows and 15{,}083{,}480
cells. We clipped 5{,}393 strictly positive cells to zero, which is 0.0358\% of the complete
matrix. Section~S2.1 reports what happens if those cells are kept instead.

\begin{table}[htbp]
\centering
\caption{The two Vina score matrices and the supports used throughout. Full SHA-256 digests,
byte counts and licences for every input are in \protect\path{results/public_core_source_manifest.csv}.}
\label{tab:s1sources}
\small
\begin{tabular}{lrr}
\toprule
 & Docking-44 & DOCKSTRING-58 \\
\midrule
Redistributed file & \path{df_final_v4.csv.gz} & \path{dockstring-dataset.tsv.gz} \\
Licence & CC BY 4.0 & Apache-2.0 \\
Source rows & 12{,}651 & 260{,}155 \\
Targets & 44 & 58 \\
Rows dropped for missingness & 0 & 95 \\
Analysis rows & 12{,}651 & 260{,}060 \\
Analysis cells & 556{,}644 & 15{,}083{,}480 \\
Missing cells in the source & 8{,}159 & 339 \\
Missing-cell rule & target-column mean & row dropped \\
Strictly positive cells & 0 & 5{,}393 \\
Positive-cell rule & clip to zero (no effect) & clip to zero \\
\bottomrule
\end{tabular}
\end{table}

Every file the evidence builder reads is declared in
\path{results/public_core_source_manifest.csv}, one row per file with its path, kind, licence,
byte count and SHA-256. The manifest has 149 rows: 111 input records, 37 code records and one
output record.
Its own digest is stored in \path{results/public_core_source_manifest.sha256};
\texttt{make verify-public-manifest} re-hashes both the manifest and every registered file.
The builder \path{analysis/build_public_core_evidence.py} opens files only through the declared
table: an undeclared key raises, a symlink anywhere along the path raises, a resolved path outside
the package raises, and 12 inputs carry a pinned digest that aborts the build on any mismatch. A
repository-wide manifest covering licence and redistribution status for every frozen file is
\path{data_manifest.csv}. The assay panels are described in S6, where they are used.

\sisubsection{S1.2 Map and dimension estimands}

Let $X$ be an $N\times P$ score matrix with ligands in rows and targets in columns. The
\emph{target map} of $X$ is the $P\times P$ sample Pearson correlation matrix of its columns.
We call its $P(P-1)/2$ strict-upper-triangle entries the \emph{edges} of the map. The
\emph{residual surface} is the two-way-centred matrix
\begin{equation}
\Gamma_{ij}=X_{ij}-\bar X_{i\cdot}-\bar X_{\cdot j}+\bar X_{\cdot\cdot},
\label{eq:s1residual}
\end{equation}
and the \emph{residual map} is the target map of $\Gamma$. Correlation already removes and
rescales column means, so the only operative difference between the two maps is subtraction of
each ligand's mean across targets. Equation~\ref{eq:s1residual} makes every residual row sum to
zero, so the residual map carries one exactly zero eigenvalue and at most $P-1$ free
directions: 43 for Docking-44 and 57 for DOCKSTRING-58.

We summarise a map by the participation-ratio dimension of its correlation eigenvalues
$\lambda_k$,
\begin{equation}
r_{\mathrm{PR}}=\frac{\left(\sum_k\lambda_k\right)^2}{\sum_k\lambda_k^2}
=\frac{P}{1+(P-1)\,\overline{r_{jk}^{\,2}}},
\label{eq:s1pr}
\end{equation}
where the bar averages the squared off-diagonal correlations \cite{roy2007,mazzucato2016}. The
second form follows from $\sum_k\lambda_k=P$ and $\sum_k\lambda_k^2=P+2\sum_{j<k}r_{jk}^2$, so
PR is an exact re-expression of one scalar: the mean squared target correlation. It runs from 1
when every $|r_{jk}|=1$ to $P$ when every $r_{jk}=0$.

What PR does not measure matters as much. It is not an algebraic rank, and no eigenvalue is
driven to zero when PR falls. Two maps with the same mean squared correlation have the same PR
however their edges are arranged, so PR by itself cannot separate one strong block of targets
from many weak pairs. PR counts no binding modes and no independent physical mechanisms, and it
is not an estimate of a molecular dimension. Where a claim needs the arrangement of edges
rather than their mean square, we report edge-rank agreement or sign agreement instead.

Every quantity here is conditional on a fixed target panel and a stated ligand support.
Docking-44 fixes 44 PDB structures selected by its source project, DOCKSTRING-58 fixes 58 gene
symbols selected by its own, and neither panel is a sample from anything. Changing $P$ changes
both the edge population and the bound on PR. Changing the library changes the correlations,
and S4 measures by how much.

\sisubsection{S1.3 Overlap of the two Vina matrices}

The producer is \path{analysis/cross_panel_overlap_audit.py} and the frozen output is
\path{results/cross_panel_overlap_audit/summary.json}. We recomputed Standard InChIKeys from
the Docking-44 analysis SMILES with the pinned RDKit \cite{rdkit2025}, and read DOCKSTRING keys
and Bemis--Murcko scaffolds \cite{murcko1996} from the frozen identity contract
\path{data/frozen/dockstring_identity_contract_2026-08-03.csv.gz}.

The two matrices come from different projects, but they are not disjoint samples
(Table~\ref{tab:s1overlap}). 582 ligands match on the full Standard InChIKey. Collapsing
stereochemical and protonation suffixes to the 14-character connectivity block gives 804 shared
blocks, covering 804 Docking-44 rows and 1{,}014 DOCKSTRING rows. 1{,}340 non-empty
Bemis--Murcko scaffolds are shared, and they cover 7{,}529 Docking-44 rows, 59.5\% of that
matrix, against 25{,}614 DOCKSTRING rows, 9.8\% of the larger one.

On the target side, a curated mapping audited against PDBe/SIFTS matches nine receptor
identities. Seven of the nine Docking-44 structures are human: 2vt4 is turkey ADRB1 (P07700)
and 3ln1 is mouse PTGS2 (Q05769), each an orthologue of the human entry DOCKSTRING scores.
Three of the nine mapped pairs dock into the same receptor structure DOCKSTRING itself used,
namely 2vt4, 3ln1 and 6cm4. The broad family labels overlap in a single entry, nuclear
receptor, out of 12 labels on the Docking-44 side and 6 on the DOCKSTRING side.

\begin{table}[htbp]
\centering
\caption{Chemical and receptor overlap between the two Vina matrices, from
\protect\path{results/cross_panel_overlap_audit/summary.json}.}
\label{tab:s1overlap}
\small
\begin{tabular}{lrr}
\toprule
Shared identity & Docking-44 rows & DOCKSTRING rows \\
\midrule
582 full Standard InChIKeys & 582 & 582 \\
804 connectivity blocks & 804 & 1{,}014 \\
1{,}340 Bemis--Murcko scaffolds & 7{,}529 (59.5\%) & 25{,}614 (9.8\%) \\
\midrule
\multicolumn{3}{l}{9 mapped receptor identities, 7 with a human Docking-44 structure} \\
\multicolumn{3}{l}{3 mapped pairs sharing the receptor structure DOCKSTRING used} \\
\multicolumn{3}{l}{1 shared broad family label (nuclear receptor) of 12 and 6} \\
\bottomrule
\end{tabular}
\end{table}

We therefore call the two matrices separately generated, and never independent. Agreement
between them is replication across matrices built by different projects with the same
underlying scoring function, not evidence from disjoint chemistry or disjoint receptors.

\sisubsection{S1.4 How to read ranges and probabilities}

No bracketed range in this Supplement is a confidence interval for a superpopulation of
targets. Both panels are fixed, so a target is not a draw from anything, and no interval here
supports a statement about proteins outside the 44 structures and 58 gene symbols analysed.
Four kinds of statement appear, and they quantify different operations.

\emph{Simulation ranges} are the 2.5th and 97.5th percentiles over repetitions of a null
construction, 500 repetitions per family in S2.3. They describe the spread of the simulation
and say nothing about sampling error in the observed value they are compared with.

\emph{Repeated-sample ranges} are the same percentiles over a stated resampling operation
applied to the observed data: 1{,}000 equal-$n$ random supports in S2.2, 200 size-matched
splits into halves sharing no whole chemical group in S4, and 100 calibration samples at each
size in S5.

\emph{Cluster bootstrap intervals} are percentile intervals from resampling whole chemical
groups with replacement, Butina clusters for Docking-44 \cite{butina1999} and Bemis--Murcko
scaffolds for DOCKSTRING, at 5{,}000 replicates per scheme. Where both schemes apply we report
the conservative union of the two intervals. The resampling unit is a chemical group and not a
ligand, so these intervals cover ligand-side variation only, with the target panel held fixed.

\emph{Permutation probabilities} are Monte Carlo target-label probabilities. One complete
relabelling of the docking target panel is applied to the rows and columns of its map while the
experimental map stays fixed \cite{hubert1976}, using 50{,}000 relabellings on the kinase and
safety-panel endpoints and 20{,}000 on the PDSP graph. All permutation and simulation
probabilities use the add-one rule, so 500 null repetitions floor at $1/501=0.0020$ and
50{,}000 relabellings floor at $1/50{,}001$.

Replicate-level values behind every range are in the frozen outputs named at the head of each
subsection, so any percentile reported here can be recomputed from the released files.
